\section{Feature engineering and preprocessing}
\label{sec:features}

\subsection{Three engineered features}

A logistic regression is linear in the log-odds: it cannot draw a curve, a
threshold or an interaction on its own. We gave it these shapes explicitly,
with one feature of each kind, each motivated by the analysis above.
\begin{itemize}
  \item \code{age\_squared} $= \text{age}^2$ (square term). With \code{age} and
  \code{age\_squared}, the model can draw the inverted U of
  Figure~\ref{fig:eda-age} instead of a straight line.
  \item \code{education\_x\_hours} $= \text{education-num} \times
  \text{hours-per-week}$ (interaction). The value of an extra hour of work is
  likely to depend on the level of education.
  \item \code{has\_capital} $= \mathbf{1}[\text{capital-gain} > 0 \text{ or }
  \text{capital-loss} > 0]$ (indicator). With more than 90\% of the amounts at
  zero, owning capital at all is a cleaner signal than the raw, skewed amounts.
\end{itemize}
The three features only combine columns of the same row. They learn nothing
from the data (no mean, no median, no target), so they cannot leak information
from validation rows into training rows. They are computed inside the
Pipeline, so they are recomputed for any new row.

\paragraph{What they bought.} On the untuned logistic regression (5-fold
out-of-fold predictions, threshold 0.5), the three features raise ROC-AUC from
0.905 to 0.913, AP from 0.763 to 0.784 and F1 from 0.657 to 0.673: about seven
times the standard deviation of AP across folds (0.003). We kept them for all
three models so that every model sees the same inputs; SHAP shows that
LightGBM still uses two of them (Section~\ref{sec:shap}).

\subsection{Missing values and encoding}

All the steps below are learned on the training rows only, inside each
cross-validation fold (Section~\ref{sec:pipeline}).
\begin{itemize}
  \item \textbf{Categorical columns} (7): missing values are replaced by the
  most frequent category, then one-hot encoded, since categories are labels,
  not measurements. \code{handle\_unknown="ignore"} turns a category never seen
  in training into a row of zeros instead of an error at prediction time.
  \item \textbf{Numerical columns} (5 original and 3 engineered): median
  imputation, then standardization $x \mapsto (x - \mu)/\sigma$, which matters
  for the regularized logistic regression. Adult has no missing numerical
  value, but the imputer stays so that a new row with a missing number can
  still be scored.
\end{itemize}
After encoding, the model sees 91 columns. Section~\ref{sec:discussion} tests a
separate ``missing'' category instead of the most frequent one.
